\documentclass[11pt]{article}
\usepackage{amsmath,amssymb,geometry,hyperref}
\geometry{margin=1in}
\newcommand{\Tr}{\operatorname{Tr}}
\begin{document}
\title{A general second-order AC-subspace bound from an ancillary spin}
\author{Exact catalogue calculation}
\date{30 October 2026}
\maketitle

\paragraph{Theorem.}
Let $N\ge8$ and let $\mathcal C\subset\operatorname{Sym}^N(\mathbb C^2)$ be a complex second-order anticoherent subspace of dimension $k$. Then
\[
\boxed{4k\le N-2,\qquad k_{\max}(N,2)\le\left\lfloor\frac{N-2}{4}\right\rfloor.}
\]
The same upper bound applies to every AC order at least two. This is an analytic all-$N$ theorem. Its proof uses no numerical optimization, reality condition on the code, or assumed symmetry of individual representatives.

\section{Whole-subspace second moments}
Write $j=N/2$, let $J_x,J_y,J_z$ be the spin-$j$ generators in units $\hbar=1$, and let $P$ project onto $\mathcal C$. The condition that every vector in $\mathcal C$ is second-order AC is equivalent to
\begin{equation}\label{eq:moments}
PJ_aP=0,\qquad PJ_aJ_bP=\frac{j(j+1)}{3}\delta_{ab}P
\qquad(a,b\in\{x,y,z\}).
\end{equation}
Indeed, the symmetric traceless part of $J_aJ_b$ has tensor rank two, its antisymmetric part is a linear combination of the rank-one spin generators, and its scalar part is $j(j+1)\delta_{ab}I/3$. Polarization gives all off-diagonal compressed matrix elements, not only the expectations of a chosen basis.

Introduce an ancillary spin $s=3/2$, with generators $S_a$ on $\mathbb C^4$, and define
\[
X=\sum_{a=x,y,z}J_a\otimes S_a,
\qquad
W=X^2-\frac{N-4}{2}X-\frac{5N(N+2)}{16}I.
\]
Let $\Pi=P\otimes I_4$. Equation~\eqref{eq:moments} gives
\[
\Pi X\Pi=0
\]
and
\begin{align*}
\Pi X^2\Pi
&=\sum_{a,b}PJ_aJ_bP\otimes S_aS_b\\
&=\frac{j(j+1)}{3}P\otimes\sum_aS_a^2\\
&=\frac{N(N+2)}{12}\frac{15}{4}\Pi
=\frac{5N(N+2)}{16}\Pi.
\end{align*}
Consequently,
\begin{equation}\label{eq:nullcompression}
\boxed{\Pi W\Pi=0.}
\end{equation}
Thus a hypothetical $k$-dimensional code gives a $4k$-dimensional subspace on which the Hermitian form defined by $W$ vanishes identically.

\section{The four eigenvalues}
The total spin $L$ in $V_{N/2}\otimes V_{3/2}$ takes the four values
\[
L=\frac{N-3}{2},\quad\frac{N-1}{2},\quad\frac{N+1}{2},\quad\frac{N+3}{2},
\]
each once. Since
\[
X=\frac12\bigl((\mathbf J+\mathbf S)^2-\mathbf J^2-\mathbf S^2\bigr),
\]
its eigenvalues are
\[
-\frac{3(N+2)}4,\quad-\frac{N+8}4,\quad\frac{N-6}4,\quad\frac{3N}4.
\]
Substitution into the quadratic polynomial defining $W$ gives
\[
\begin{array}{c|c|c}
L&\text{eigenvalue of }W&\text{multiplicity}\\\hline
(N-3)/2&(N+2)(5N-3)/8&N-2\\
(N-1)/2&-N(N-7)/8&N\\
(N+1)/2&-(3N^2+N+6)/8&N+2\\
(N+3)/2&-N(N-7)/8&N+4
\end{array}
\]
For every $N\ge8$ the first eigenvalue is strictly positive and the other three are strictly negative. There are no zero eigenvalues. Hence the inertia of $W$ is
\[
(n_+,n_-,n_0)=(N-2,3N+6,0).
\]
These signs follow directly from the displayed factors for all $N$ in the theorem's domain; checking finitely many spectra is unnecessary.

\section{The dimension argument}
Let $V_+$ and $V_-$ be the positive and negative spectral subspaces of $W$. If a subspace $\mathcal U$ satisfies $\langle u,Wv\rangle=0$ for every $u,v\in\mathcal U$, its orthogonal projection onto $V_+$ is injective: a nonzero vector in the projection kernel would lie entirely in $V_-$ and would satisfy $\langle u,Wu\rangle<0$, a contradiction. Therefore
\[
\dim\mathcal U\le\dim V_+=N-2.
\]
Apply this to $\mathcal U=\mathcal C\otimes\mathbb C^4$ using~\eqref{eq:nullcompression}. Its dimension is $4k$, proving the theorem.

\paragraph{Why seven qubits are not excluded.}
At $N=7$, the second and fourth eigenvalues of $W$ vanish, with combined multiplicity $18$. The injectivity argument using the strictly negative complement then fails. This is consistent with the known exact two-dimensional second-order AC subspace at $N=7$; the theorem deliberately assumes $N\ge8$.

\section{The thirteen-qubit case}
At $N=13$, the canonical witness is
\[
W=X^2-\frac92X-\frac{975}{16}I.
\]
Its eigenvalues and multiplicities are
\[
\frac{465}{4}\ (11),\qquad-\frac{39}{4}\ (13),\qquad
-\frac{263}{4}\ (15),\qquad-\frac{39}{4}\ (17).
\]
Thus $4k\le11$, excluding dimension three. An exact two-dimensional witness, written in normalized Dicke states, is
\[
|u\rangle=\frac{\sqrt{10}|D_{13}^{(0)}\rangle+\sqrt{26}|D_{13}^{(9)}\rangle}{6},
\qquad
|v\rangle=\frac{-\sqrt{26}|D_{13}^{(4)}\rangle+\sqrt{10}|D_{13}^{(13)}\rangle}{6}.
\]
These vectors are orthonormal and obey all the identities~\eqref{eq:moments}, including the cross terms. Hence
\[
\boxed{k_{\max}(13,2)=2.}
\]
This is a global exclusion: it applies to every possible three-dimensional complex subspace and does not merely rule out extensions of the displayed plane.

\paragraph{Further consequence.}
At $N=16$ and $N=17$, the bound gives $k\le3$. The catalogue already contains exactly verified three-dimensional second-order witnesses at both sizes, so these maxima are also three. More generally, the formula can be combined with any independently certified lower bound.

\paragraph{Reproducibility and scope.}
The accompanying \texttt{verify\_general.py} verifies the four eigenvalue formulas as symbolic identities in $N$, the signs for the entire range $N\ge8$, and the ancillary Casimir. It also constructs the full physical witness matrices and checks all compressed entries exactly for the displayed $N=13$ plane and the known $N=7$ code. The independent argument is the analytic inertia proof above; finite checks validate conventions and implementation. No claim of literature priority is made.
\end{document}
